% Folder 134-2, batch 03 — archivists' pages 41–60.
% Pass: Opus 5.5 (claude-opus-5-5), 2026-10-03 — first pass, unchecked against the pages by a human.
%
% Special rules for folder 134-2 (issue #44) applied: Pursuing Stacks text is not
% transcribed; Grothendieck's 1975 letters to Lawrence Breen are; third-party
% letters are not. This batch holds no Pursuing Stacks text and no third-party
% letter: all twenty pages are transcribed.
%
% Letters in this batch:
%  - pp. 41–47 (his pp. 26–32): end of a typed copy of a French letter in his
%    hand, annotated by him; opening, date and recipient fall before the batch
%    and are not visible here. Signed « A Grothendieck », P.S., « Re-salut … A.G. ».
%    Very probably the long 1975 letter to Breen: p. 48 calls itself an
%    « afterthought à une lettre-fleuve sur le yoga homotopique » and p. 50 refers
%    to the fear of « n-topos » « exprimée dans ma précédente lettre », which is
%    p. 46. Inference only; the pages do not name the recipient. Transcribed.
%  - pp. 48–55 (his pp. 33–39 bis, letter pp. 1–8): « Villecun le 17.2.1975 »,
%    to Lawrence Breen (« Cher Larry »), French, typed copy with his handwritten
%    script letters and corrections; struck heading « Transcription d'une
%    lettre manuscrite de A. Grothendieck ». Transcribed in full.
%  - pp. 56–60 (his pp. 40, 42–45; letter p. 1 then « -12- » to « -15- »):
%    « Villecun 17/19 July 1975 », to Lawrence Breen (« Dear Larry »), English,
%    typed with his handwritten date and corrections. Letter pp. 2–11 are not in
%    this batch; the letter runs on past p. 60. Transcribed.
% Pages where it is unclear whether the text belongs to Pursuing Stacks: none.
% The marginal « (App. N) » marks (3 to 11) and footnote calls (25)–(27) are in
% his hand, apparently from a later reworking of these letters (1983?) — this is
% an inference, not on the page.
\documentclass[11pt,a4paper]{article}
\input{../preamble/grothendieck}

\folder{134-2}
\batch{3}
\pages{41}{60}
\foldertitle{[Chapitre I : Take off (pages 1 à 65) et table des matières provisoire] : tapuscrits annotés (19/02-22/02/1983), note manuscrite (s.d.), copies de lettre (1975, s.d.).}
\dating{1975-[1983]}
\watermark{Édition de démonstration}

\begin{document}

\section{Lettre (fin), destinataire et date hors du lot}

\page{41}
\note{p. 26 de l'auteur. Copie dactylographiée d'une lettre de Grothendieck, annotée de sa main ; le début de la lettre est antérieur à ce lot, et la phrase reprend en cours.}
\marginal{$\pi_0$, $\pi_1$}
… à coéfficients discrets\add{ en cohomologie étale}, dans le cas d'une variété projective disons et de toute section hyperplane, ou d'une variété quasi-projective et de presque toute section hyperplane (pour ne mentionner que le cas global le plus simple), sous les hypothèses de profondeur \add{cohomologique} « les plus naturelles » (en fait, essentiellement des conditions néc. et suff. de validité dudit théorème). Dans le cas commutatif, les techniques de dualité nous suggèrent très clairement quels sont les meilleurs énoncés possibles, cf. l'exposé de Mme Raynaus\note{sic, pour « Raynaud »} dans SGA~2. Mais ces techniques ne valent qu'en se restreignant à des coéfficients premiers aux caractéristiques, alors que des démonstrations directes plus géométriques (développées dans SGA~2 avant le développement du formalisme de la cohomologie étale) donnaient des résultats très voisins pour le $H^0$ et le $H^1$ (ou le \add{$\pi_0$ et le} $\pi_1$, si on préfère) sans telles restrictions, du moins dans le cas propre (i.e. projectif, au lieu de quasi-projectif). En fait, ce sont \struck{en fait} les « résultats les meilleurs possibles » eux-mêmes, énoncés comme conjectures dans SGA~2 dans l'exposé cité de Mme Raynaud, qui sont démontrés ultérieurement par elle dans sa thèse. Ce qui est remarquable de notre point de vue, c'est que les énoncés les plus forts se présentent le plus naturellement sous forme d'énoncés sur des \add{1-}\emph{champs} sur le site étale de la variété algébrique considérée -- la notion de « profondeur $\geqslant i$ » (pour $i = 1, 2, 3$) s'énonçant aussi le plus naturellement en termes de champs. Non seulement cela, mais alors même qu'on voudrait ignorer la notion technique de champ et travailler exclusivement en termes de $H^0$ et $H^1$ en utilisant à bloc le formalisme \add{cohomologique non commutatif} de Giraud, pour démontrer disons un théorème de bijectivité $\pi_1(Y) \to \pi_1(X)$ (ce qui est le résultat le plus profond établi dans la thèse de Mme Raynaud), il semble bien qu'on n'y arrive pas, faute à ce formalisme d'avoir la souplesse nécesssaire. En fait, il faut utiliser comme ingrédients techniques, de façon essentielle, les trois théorèmes suivants directement pour les 1-champs « de torsion » (i.e. où les faisceaux en groupes d'automorphismes sont de ind-torsion) : a) théorème de changement de base pour un morphisme propre, b) théorème de changement de base par un morphisme lisse c) théorème de \struck{chang} « propreté cohomologique générique » pour un morphisme de type fini $f\colon X \to S$, $S$ intègre (disant que l'on peut trouver dans $S$ un ouvert $U \neq \emptyset$ tel que pour \emph{tout} changement de base $S' \to S$ se factorisant par $u$, la formule de changement de base est vraie). (Pour b) et c), il faut faire des hypothèses que les faisdeaux d'automorphismes sont premiers aux carafctérisitques, et dans c) une hypothèse de constructtibilité sur ces faisceaux -- mais b) et c) ne servent que dans la version « générique » du th. de Lefschetz.) C'est avec en vue de telles applications que Giraud a pris la peine dans son bouquin (si je ne me trompe) de demontrer a), b) (et c) ?) dans le contexte des 1-champs et de leurs

\page{42}
\note{p. 27 de l'auteur.}
images directes et inverses. Mais du même coup il devient clair que le contexte « naturel » des théorèmes de changement de base en cohomologie étale, des théorèmes \struck{de} \add{du type} Lefschetz (dits « faibles ») sur les « sections hyperplanes », tout comme \add{de} la notion de profondeur qui y joue un rôle crucial, \add{doit être} celui des $n$-champs. Et que le développement hypothétique de ce contexte, ne risque pas de se réduire à une jonglerie purement formelle et absolument bordélique avec du « general nonsense », mais qu'on se trouvera aussitôt confronté à des tests « d'utilisabilité » aussi sérieux que la démonstration des théorèmes de changement de base et ceux du type de Lefschetz (qui même dans le contexte commutatif ne sont pas piqués de vers …). \add{\uncertain{Idem} pour variantes analytiques complexes etc.}

Je ne sais si ces commentaires te « passent par dessus la tête » à ton tour, ni si elles te donnent l'impression qu'il y aurait peut-être des choses intéressantes à tirer au clair. Si cela t'intéresse, je pourrais expliciter sous forme un peu plus systématique quelques ingrédients d'une hypothétique algèbre homologique non commutative et les liens de celle-ci à l'algèbre homologique commutative. Plus mystérieux pour moi (et pour cause, vu mon ignorance en homotopie) seraient les relations entre celle-là et l'algèbre homotopique, i.e. les structures semi-simpliciales, et je n'ai que des commentaires assez vagues à faire en ce sens\add{(*)}. Par ailleurs, je te rappelle que même l'algèbre homologique commutative n'est pas, il s'en faut, dans un état satisfaisant, pour autant que je sache, vu qu'on ne sait toujours pas quelle est la « bonne » notion de catégorie triangulée\add{(4)}. Or il me semble bien clair que ce n'est pas une question purement académique -- même si on a pu se passer de le savoir jusqu'ici (en se bornant comme Monsieur Jourdain à « faire de la prose sans le savoir » -- en travaillant sur des catégories de complexes, éventuellement filtrés, sans trop se demander quelles structures il y a sur ces catégories …).
\marginal{(\uncertain{24})}
\note{les appels (*) et (4) sont ajoutés à la main ; le (*) renvoie au P.S. ci-dessous.}

Bien cordialement à toi

A Grothendieck

\marginal{(App.~3)}
\add{(*) P.S.} Réflexion faite, j'ai quand même envie de te mettre un peu en appétit, en faisant ces « quelques commentaires assez vagues ». Il s'agit du yoga qu'une \add{(\ill{} (petite))} $n$-catégorie \add{(en groupoïdes)} à $n$-équivalence près \struck{\ill{}} \add{« est} essentiellement la même chose » qu'un ensemble semi-simplicial pris à homotopie près et où on néglige les $\pi_i$ pour $i \geqslant n+1$ (où, si tu préfères, « où on a tué les groupes d'homotopie en dimension $\geqslant n+1$ »). Voici des éléments heuristiques pour ce yoga. Si $K_\bullet$ est un ensemble simplicial (il peut être prudent de le prendre de Kan) on lui associe une $n$-catégorie $C_n(K_\bullet)$, dont les 0-objets sont les
\note{il écrit « $K.$ » ; transcrit $K_\bullet$ dans tout le lot.}

\page{43}
\note{p. 28 de l'auteur.}
0-simplexes, les 1-objets sont les chemins (ou homotopies) entre 0-simplexes, les 2-objets sont les homotopies entre chemins (à extrémités fixées) etc. Pour les $n$-objets, cependant, on ne prend pas les homotopies entre homotopies de fourbis, mais classes d'équivalence \struck{\ill{}} homotopies (modulo la relation d'homotopie) entre homotopies. La composition des $i$-objets ($i \geqslant 1$) se définit de façon évidente, on notera qu'elle n'est pas strictement associative, mais associative modulo homotopie. Donc la $n$-catégorie qu'on obtient n'est pas « stricte » -- et on prévoit pas mal d'emmerdements pour définir de façon raisonnable une $n$-catégorie pas stricte (dans la description des compatibilités pour les « données d'associativité »). La mise sur pied du yoga qui suit pourrait constituer un fil d'Ariane pour la définition en forme des $n$-catégories (pas strictes), les $n$-foncteurs entre elles (pas non plus stricts, et pour cause), les $n$-équivalences etc, au même titre que le yoga initial « une $n$-catégorie est une catégorie où les Hom et leurs accouplements \add{de composition} sont des $(n-1)$-catégories et des accouplements entre telles ». Cette $n$-catégorie \add{$C_n(K_\bullet)$} dépend fonctoriellement de $K_\bullet$, tout morphisme simplicial $K_\bullet \to K'_\bullet$ définit un $n$-foncteur $C_n(K_\bullet) \to C_n(K'_\bullet)$ ; en fait, cela doit en dépendre même $n$-fonctoriellement, vu qu'on voit (en s'inspirant de ce qui précède et l'appliquant à des ensembles semi-simpliciaux de la forme $\mathrm{Hom}_\bullet(K_\bullet, K'_\bullet)$) que les ensembles semi-simpliciaux forment eux-mêmes \add{les 0-objets d'une} $n$-catégorie, quel que soit $N$ …
\note{il écrit « $K!$ » pour $K'_\bullet$ (apostrophe frappée sur le point), et « quel que soit $N$ » après « $n$-catégorie ».}

En fait, $C_n(K_\bullet)$ est un $n$-groupoïde, i.e. une $n$-catégorie où toute $i$-flèche ($1 \leqslant i \leqslant n$) ($=$ $i$-objet) est une « équivalence » i.e. admet un quasi-inverse (donc un inverse si la $n$-catégorie est « réduite »). Si $C$ est une telle $n$-catégorie i.e. un $n$-groupoïde, et $x$ un 0-objet de $C$, il s'impose de désigner par $\pi_i(C, x)$ ($0 \leqslant i \leqslant n$) successivement : l'ensemble des classes de 0-objets à équivalence près (c'est un vulgaire ensemble), l'ensemble des classes à équivalence près de 1-objets \add{(ou 1-flèches)} $x \to x$ (c'est un groupe, pas nécessairement commutatif), l'ensemble des \add{classes mod équivalence des} 2-flèches $1_x \to 1_x$, où $1_x$ est la 1-flèche identique de $x$ (c'est un groupe commutatif $\pi_2(C, x)$, ainsi que les groupes qui vont suivre), l'ensemble des classes mod équivalence de 3-flèches $1_{1_x} \to 1_{1_x}$, etc. Ces groupes forment, comme de juste, des « systèmes locaux » sur l'ensemble des 0-objets de $C$, et modulo le grain de sel habituel, \add{les $\pi_i(C, x)$} ne dépendent que de la « composante connexe » du 0-objet $x$ i.e. de sa classe modulo équivalence de 0-objets. Ceci dit, si $C$ est de la forme $C_n(K_\bullet)$, il résulte pratiquement des définitions que l'on a des isomorphismes canoniques $\pi_i(K_\bullet, x) \simeq \pi_i(C_n(K_\bullet))$ pour $0 \leqslant i \leqslant n$, qui pour $x$ variable peuvent s'interpréter comme des isomorphismes de systèmes locaux. Il s'ensuit que pour une application semi-simpliciale $f\colon K_\bullet \to K'_\bullet$, le $n$-foncteur correspondant $C_n(K_\bullet) \to C_n(K'_\bullet)$ est une $n$-équivalence sss $f$ induit un iso-

\page{44}
\note{p. 29 de l'auteur.}
morphisme sur les $\pi_0$ et sur les $\pi_i$ en tout point ($1 \leqslant i \leqslant n$). On serait plus heureux de pouvoir dire à la place « et \add{de plus,} un homomorphisme surjectif pour $i = n+1$ », car c'est, il me semble, cela qu'il faudrait pour \add{espérer} pouvoir conclure que la catégorie localisée de la catégorie des ensembles semi-simpliciaux, obtenue en inversant les flèches « qui induisent des isomorphismes sur les $\pi_i$ pour $0 \leqslant i \leqslant n$ » (ou encore, « en négligeant » les ensembles semi-simpliciaux $n$-connexes), est équivalente à la catégorie localisée de la catégorie des $n$-catégories, où on rend inversibles les $n$-équivalences ? Quoi qu'il en soit, ces petites bavures devraient disparaitre lorsqu'on « stabilise » en faisant augmenter $n$. A ce propos, on voit que le foncteur « troncature en dim.~$n$ » de la théorie homotopique (consistant à tuer les groupes d'homotopie à partir de la dimension $n+1$) s'interprète dans le langage des $n$-catégories par l'opération faisant passer d'une $N$-catégorie ($N > n$) à une $n$-catégorie, en conservant tels quels les $i$-objets ($0 \leqslant i \leqslant n-1$) et leur composition ($1 \leqslant i \leqslant n-1$), et en remplaçant les $n$-objets par les classes de $n$-objets « à équivalence près », avec la composition obtenue par passage au quotient. De même, le foncteur d'inclusion évident en théorie homotopique, consistant à regarder un ensemble semi-simplicial « où on a négligé les $\pi_i$ pour $i \geqslant n+1$ » comme un ensemble semi-simplicial (dans la catégorie homotopique) qui se trouve avoir des $\pi_i$ nuls pour $i \geqslant n+1$, se traduit par le foncteur allant des $n$-catégories vers les $N$-catégories, obtenue en ajoutant à une $n$-catégorie des $i$-flèches ($n+1 \leqslant i \leqslant N$) identiques exclusivement. (Ainsi, un ensemble est regardé comme une catégorie « discrète », une catégorie comme une 2-catégorie où les $\underline{\mathrm{Hom}}(A, B)$, $A$ et $B$ des 0-objets, sont des catégories discrètes, etc …).
\note{« $i = n+1$ » : la frappe porte « $N+1$ », le $N$ surchargé en $n$.}

Bien entendu, rien n'empêche de considérer aussi la notion de $\infty$-catégorie, à laquelle celle de $n$-catégorie est comme la notion d'ensemble semi-simplicial tronqué à celle d'ensemble semi-simplicial. Sauf erreur, la localisée de la catégorie des $\infty$-catégories, pour les flèches de $\infty$-équivalence, est \add{équivalente à} « la catégorie homotopique », localisée de la catégorie des ensembles semi-simpliciaux, ou du moins une sorte de complétée de celle-là. Dans cette optique, le tapis consistant à interpréter une $\infty$-catégorie de Picard stricte (i.e. quelque chose qui ressemble à un groupe abélien de la catégorie des $\infty$-catégories) comme donnée (à $\infty$-équivalence près) par un complexe de chaines regardé comme un objet d'une catégorie dérivée, est à relier au tapis de Dold-Puppe, interprétant ces derniers comme des groupes abéliens semi-simpliciaux.

Pour se donner confiance dans ce yoga \add{général}, on peut essayer d'interpéter en termes de $n$-catégories ou $\infty$-catégories des constructions familières

\page{45}
\note{p. 30 de l'auteur.}
en homotopie. Ainsi, l'espace des lacets $\Omega(K_\bullet, x)$ correspond manifestement à la $(n-1)$-catégorie $\underline{\mathrm{Hom}}(x, x)$ formée des $i$-flèches de $C$ ($1 \leqslant i \leqslant n$) dont la 0-origine et la 0-extrémité sont $x$, réindexées en les appelant $(i-1)$-flèches. Je n'aperçois pas à vue de nez un joli candidat pour la suspension en termes de $n$-catégories. Par contre le $\mathrm{Hom}_\bullet(K_\bullet, K'_\bullet)$ doit correspondre au $\underline{\mathrm{Hom}}(C, C')$, qui est une $n$-catégorie quand $C, C'$ en sont. La « fibre homotopique » d'une application semi-simpliciale $f\colon K_\bullet \to K'_\bullet$ (transformée d'abord, pour les besoins de la cause, en une fibration de Serre par le procédé bien connu de Serre-Cartan) correspond sans doute à l'opération bien familière de produit $n$-fibré (du moins les cas $n = 0, 1$ sont bien familiers !) $C \times_{C'} C''$ pour des $n$-foncteurs $C \to C'$ et $C'' \to C'$, dans le cas où $C''$ est la $n$-catégorie ponctuelle, donc la donnée de $C'' \to C'$ correspond à la donnée d'un 0-objet de $C'$. Les espaces $K(\pi, n)$ ont une interprétation évidente comme $n$-gerbes liées par $\pi$. Enfin, on voit aussi poindre l'analogue du dévissage de Postnikoff d'un ensemble semi-simplicial -- mais la façon dont je l'entrevois (vue ma prédilection pour les topos) passe par la notion de topos classifiant d'une $n$-groupoïde (généralisant de façon évidente le topos classifiant d'un groupe). En termes de cette notion, on peut, il me semble, interpréter un $n$-groupoïde en termes d'un $(n-1)$-groupoïde (savoir son tronqué), \emph{muni} d'une $n$-gerbe sur le topos classifiant, liée par $\pi_n$ \add{(« tordu » bien sûr par l'action du $\pi_1$ …)}.
\note{« produit $n$-fibré » : une main ajoute « $+1$ » et une parenthèse autour du $n$, et écrit « $(n+1)$ » dans l'interligne, ce qui corrige vraisemblablement en « produit $(n+1)$-fibré » ; la portée exacte de la correction n'est pas sûre.}

\marginal{\uncertain{espacer}}
\marginal{(App.~4)}
Bien sûr, il faut relativiser encore tout le yoga qu'on vient de décrire, au dessus d'un topos quelconque $X$. Il s'agirait donc de mettre en relation et d'identifier, dans une certaine mesure, d'une part l'algèbre homotopique sur $X$ construite en termes de la notion de faisceau semi-simplicial sur $X$, d'autre part l'algèbre catégorique sur $X$ construite en termes de la notion de $n$-champ en groupoïdes ($n \geqslant 0$ fini ou infini). On espère que la notion d'image inverse de faisceau semi-simplicial par un morphisme de topos $f\colon X \to X'$ (qui est évidente) correspond à la notion \add{plus} subtile d'image inverse de $n$-champs ; et inversement, la notion évidente d'image directe de $n$-champs par $f$ devrait correspondre à une notion plus subtile d'image directe \add{$Lf_*(K_\bullet)$} d'un faisceau semi-simplicial, construit sans doute dans l'esprit des foncteurs dérivés à partir de la notion naive (mais on hésite s'il faut mettre $Lf_*$ ou $Rf_*$) … Les dévissages à la Postnikoff doivent avoir encore une interprétation remarquablement simple en termes de $n$-champs. Comparer à la remarque de Giraud qu'un (1)-champ en groupoides sur $X$ peut s'identifier au couple d'un faisceau $\pi_0$ sur $X$, et d'une (1)-gerbe sur le topos induit $X_{/\pi_0}$ (dont le lien, comme de juste, devrait être noté $\pi_1$ !). D'ailleurs, dans le cas des 1-champs en groupoides, la traduction de ces animaux en termes de topos
\note{« $Lf_*(K_\bullet)$ » est écrit à la main au-dessus de « d'un faisceau », qu'un trait semble barrer en partie ; on ne sait s'il remplace ces mots ou les précise.}

\page{46}
\note{p. 31 de l'auteur.}
classifiants au dessus de $X$ est, je crois, développé en long et en large dans Giraud (il parle, si je me rappelle bien, d'« extensions » du topos $X$). L'extension (si j'ose dire) de ce tapis aux $n$-champs ne devrait pas poser de problème.
\note{après « Giraud » la frappe porte un « x » en indice, peut-être un appel de note.}

Remords : tâchant de préciser heuristiquement la notion de topos classifiant d'un $n$-champ en groupoides (ou plus particulièrement, d'un $n$-groupoide), pour $n \geqslant 2$, je vois que je n'y arrive pas à vue de nez. (Bien sûr, il suffirait (procédant de proche en proche) de savoir définir un topos classifiant raisonnable pour une $n$-\emph{gerbe}, liée par un faisceau abélien $\pi_n$.) Donc je ne sais comment décrire le dévissage de Postnikoff en termes de $n$-champs, sauf pour $n \leqslant 2$. Ceci est lié à la question d'une description directe des groupes de cohomologie d'un $n$-groupoide \add{$C$} (ou d'un $n$-champ), à coéfficients disons dans un système local commutatif, de façon que pour $C = C_n(K_\bullet)$, $K_\bullet$ un ensemble semi-simplicial dont les $\pi_i$ pour $i \geqslant n+1$ sont nuls, on trouve les groupes de cohomologie correspondants de $K_\bullet$. Peut-on le faire en associant à $C$, de façon convenable, un ensemble semi-simplicial « nerf » de $C$ ? Bien entendu, si on réussit à définir un topos classifiant pour $C$, celui-ci devrait être homotope à $K_\bullet$ ci-dessus, donc avoir les mêmes invariants homotopiques $\pi_i$ et cohomologiques $H^i$ ; itou pour les champs. La définition habituelle du topos classifiant, dans le cas $n = 1$, a bien cette vertu. Cas particulier typique du Pb de la définition du topos classifiant : pour $\pi$ un groupe commutatif, trouver un topos canonique (fonctoriel en $\pi$ bien sûr …) ayant le type d'homotopie de $K(\pi, n)$, et qui généralise la définition du topos classifiant pour $n = 1$ (topos des ensembles où $\pi$ opère). On frémit à l'idée que les topos pourraient ne pas faire l'affaire, et qu'il y faille des « $n$-topos » !! (J'espère bien que ces animaux n'existent pas …) ( ).

\marginal{\uncertain{espacer}}
\marginal{(App.~5)}
La théorie d'« algèbre homologique non commutative » que j'essaie de suggérer pourrait se définir, vaguement, comme l'étude parallèle des notions suivantes et de leurs relations multiples : a) espaces topologiques, topos, b) ensembles semi-simpliciaux, faisceaux semi-simpliciaux etc. c) $n$-catégories (notamment $n$-groupoïdes), $n$-champs (notamment $n$-champs en groupoides) etc. d) complexes de groupes abéliens, de faisceaux abéliens etc. (Les « etc » réfèrent surtout aux structures supplémentaires qu'on peut envisager sur les objets du type envisagé …) C'est donc de l'algèbre avec la présence constante de motivations provenant de l'intuition topologique. Si une telle théorie devait voir le jour, il lui faudrait bien un nom, je me demande si « algèbre topologique » ne serait pas le plus adéquat (« algèbre homologique non commutative » ne peut guère aller à la longue, pour des raisons évidentes). Ce qui est aujourd'hui parfois désigné sous ce \add{vocable} n'est guère qu'un bric à brac de notions (telles que anneau topologique

\page{47}
\note{p. 32 de l'auteur.}
corps topologique, groupe topologique \add{etc)}, \struck{qui} ne forme\struck{nt} guère un corps de doctrine cohérent -- il ne s'impose donc pas que cela accapare un nom qui servirait mieux d'autres usages. (Comparer le nouvel usage du terme « géométrie analytique » introduit par Serre, et qui ne semble guère avoir rencontré de résistance.)
\note{la correction manuscrite de cette phrase (« etc) », « qui » et la fin de « forment » barrés) est lue avec doute.}

Re-salut, et au plaisir de te lire

A.G.
\note{fin de la lettre commencée avant ce lot (pp. 26--32 de l'auteur ici).}

\section{Lettre à Lawrence Breen, Villecun, 17.2.1975}

\page{48}
\note{p. 33 de l'auteur. Copie dactylographiée d'une lettre manuscrite de Grothendieck à Lawrence Breen ; l'intitulé dactylographié est barré à la main. Les lettres cursives $\mathcal{F}$ (champs) sont tracées à la main dans les blancs de la frappe.}
\struck{Transcription d'une lettre manuscrite de A. Grothendieck}

Villecun le 17.2.1975

Cher Larry,

\marginal{(App.~6)}
Encore un « afterthought » à une lettre-fleuve sur le yoga homotopique. Comme tu sais sans doute, à un topos $X$ on associe canoniquement un pro-ensemble simplicial, donc un « pro-type d'homotopie » en un sens convenable. Dans le cas où $X$ est « localement homotopiquement trivial », le pro-objet associé est essentiellement constant en tant que pro-objet dans la catégorie homotopique, donc $X$ définit un objet de la catégorie homotopique usuelle, qui est son « type d'homotopie ». De même, si $X$ est « localement homotopiquement trivial en dim~$\leqslant n$ », il définit un type d'homotopie ordinaire « tronqué en dim~$\leqslant n$ » -- construction familière pour $i = 0$ ou $1$, même à des gens comme moi qui ne connaissent guère l'homotopie !

Ces constructions sont fonctorielles en $X$. D'ailleurs, si $f\colon X \to Y$ est un morphisme de topos, Artin-Mazur ont donné une condition nécéssaire et suffisante \emph{cohomologique} pour que ce soit une « équivalence d'homotopie en dim~$\leqslant n$ » : c'est que $H^i(Y, F) \xrightarrow{\sim} H^i(X, f^*(F))$ pour $i \leqslant n$, et tout faisceau de groupes \emph{localement constant} $F$ sur $Y$, en se restreignant de plus à $i \leqslant 1$ dans le cas $F$ non commutatif. Ce critère, en termes de $n$-gerbes « localement constantes » $\mathcal{F}$ sur $Y$, s'interprète par la condition que $\mathcal{F}(Y) \to \mathcal{F}(X)$ est une $n$-équivalence pour tout tel $\mathcal{F}$ et $i \leqslant n$. Il est certainement vrai que ceci équivaut encore au critère suivant
\note{la flèche $\mathcal{F}(Y) \to \mathcal{F}(X)$ est tracée à la main, avec une boucle à l'origine (peut-être un $\sim$).}

\begin{quote}
(A) Pour tout $n$-champ « localement constant » $\mathcal{F}$ sur $Y$, le $n$-foncteur $\mathcal{F}(Y) \to f^*(\mathcal{F})(X)$ est une $n$-équivalence
\end{quote}

ou encore à

\begin{quote}
(B) Le $n$-foncteur $\mathcal{F} \mapsto f^*(\mathcal{F})$ allant de la $n$-catégorie des $(n-1)$-champs localement constants sur $Y$ dans celle des $(n-1)$-champs localement constants sur $X$, est une $n$-équivalence.
\end{quote}

En d'autres termes, les constructions sur un topos $X$ qu'on peut faire en termes de $(n-1)$ champs \emph{localement} constants ne dépendent que de

\page{49}
\note{p. 34 de l'auteur ; la frappe porte aussi « 2 », pagination propre de la lettre.}
son « (pro)-type d'homotopie $n$-tronquée », et le définissent. Dans le cas où $X$ est localement homotopiquement trivial en dim~$\leqslant n$, donc définit un type d'homotopie $n$-tronqué ordinaire, on peut interpréter ce dernier comme un $n$-groupoide $\mathcal{C}_n$, (défini à $n$-équivalence près). En termes de \struck{\ill{}} \add{$\mathcal{C}_n$,} les $(n-1)$-champs \add{localement constants} sur $X$ doivent s'identifier aux $n$-foncteurs de la $n$-catégorie $\mathcal{C}_n$ dans la $n$-catégorie $((n-1)\text{-}\mathrm{Cat})$ de toutes les $(n-1)$-catégories. Dans le cas $n = 1$ ceci n'est autre que la théorie de Poincaré de la classification des revêtements de $X$ en termes du « groupoide fondamental » $\mathcal{C}_1$ de $X$. Par extension, $\mathcal{C}_n$ mérite le nom de \emph{$n$-groupoide fondamental de $X$}, que je propose de noter $\Pi_n(X)$. Sa connaissance induit donc celle des $\pi_i(X)$ ($0 \leqslant i \leqslant n$) et des invariants de Postnikoff de tous les ordres jusqu'à $H^{n+1}(\Pi_{n-1}(X), \pi_n)$.
\note{dans la marge gauche, un crochet hachuré marqué « (C) » encadre le passage « En termes de $\mathcal{C}_n$, … de toutes les $(n-1)$-catégories. », que deux équerres délimitent dans le texte.}

Dans le cas d'un topos $X$ quelconque, pas nécéssairement \add{loc.} homotopiquement trivial en dim~$\leqslant n$, on espère pouvoir interpréter les $(n-1)$-champs localement constants sur $X$ en termes d'un $\Pi_n(X)$ qui sera un pro-$n$-groupoide. Ça a été fait en tous cas, plus ou moins, pour $n = 1$ (du moins pour $X$ connexe) ; le cas où $X$ est le topos étale d'un schéma est traité in extenso dans SGA~3, à propos de la classification des tores sur une base quelconque.

Dans le cas $n = 1$, on sait qu'on récupère (à équivalence près) le 1-groupoide $\mathcal{C}_1$ à partir de la 1-catégorie $\underline{\mathrm{Hom}}(\mathcal{C}_1, (\mathrm{Ens}))$ de ces foncteurs dans $(\mathrm{Ens}) = (0\text{-}\mathrm{cat})$ (i.e. des « systèmes locaux » sur $\mathcal{C}_1$) (qui est un topos, dit « multigaloisien ») comme la catégorie des « foncteurs fibres » sur ledit topos, i.e. la catégorie opposée à la catégorie des points de ce topos (\struck{\ill{}} \add{lequel} n'est autre que le \emph{topos classifiant} de $\mathcal{C}_1$). Pour préciser pour $n$ quelconque la façon dont le $n$-type d'homotopie d'un topos $X$ (supposé localement homotopiquement trivial en dim~$\leqslant n$, pour simplifier), i.e. son $n$-groupoide fondamental $\mathcal{C}_n$, s'exprime en termes de la $n$-catégorie des

\page{50}
\note{p. 35 de l'auteur ; pagination de la lettre : 3.}
« $(n-1)$-systèmes locaux sur $X$ » i.e. des $(n-1)$-champs localement constants sur $X$, et par là élucider complètement l'énoncé hypothétique (B) ci-dessus, il faudrait donc expliciter comment un $n$-groupoide $\mathcal{C}_n$ se récupère, à $n$-équivalence près, par la connaissance de la $n$-catégorie
\[
\mathcal{T}_n = (n\text{-}\underline{\mathrm{Hom}})(\mathcal{C}_n, ((n-1)\text{-}\mathrm{Cat}))
\]
des $(n-1)$-systèmes locaux sur $\mathcal{C}_n$. On aurait envie de dire que $\mathcal{C}_n$ est la catégorie des « $n$-foncteurs fibres » sur $\mathcal{T}_n$, i.e. des $n$-foncteurs $\mathcal{T}_n \to ((n-1)\text{-}\mathrm{Cat})$ ayant certaines propriétés d'exactitude (pour $n = 1$, c'était la condition d'être les foncteurs image inverse pour un morphisme de topos, i.e. de commuter aux $\varinjlim$ quelconques et aux $\varprojlim$ finies …) C'est ici que se matérialise la peur, exprimée dans ma précédente lettre, qu'on finisse par tomber sur la notion de $n$-topos et morphismes de tels ! $\mathcal{T}_n$ serait un $n$-topos (appelé le « $n$-topos classifiant du $n$-groupoide $\mathcal{C}_n$ »), $((n-1)\text{-}\mathrm{Cat})$ serait le $n$-topos « ponctuel » type, et $\mathcal{C}_n$ s'interprète mod $n$-équivalence comme la $n$-catégorie \add{des} « $n$-points » du $n$-topos classifiant $\mathcal{T}_n$. Brr !
\note{les lettres cursives $\mathcal{C}_n$, $\mathcal{T}_n$, $\mathcal{B}$ sont tracées à la main ; le $\mathcal{C}_n$ et le $\mathcal{T}_n$ se ressemblent et ont été départagés par le sens.}

Si on espère encore pouvoir définir un bon vieux (1)-topos classifiant pour un $n$-groupoide $\mathcal{C}_n$, comme solution d'un Pb universel, je ne vois guère que le Pb universel suivant : pour tout topos $T$, considérons $\underline{\mathrm{Hom}}(\Pi_n(T), \mathcal{C}_n)$. C'est une $n$-catégorie, mais prenons-en la 1-catégorie tronquée $\tau_1 \underline{\mathrm{Hom}}(\Pi_n(T), \mathcal{C}_n)$. Pour $T$ variable, on voudrait 2-représenter le 2-foncteur contravariant $(\mathrm{Top})^\circ \to (1\text{-}\mathrm{Cat})$ par un topos classifiant $\mathcal{B} = \mathcal{B}_{\mathcal{C}_n}$, donc trouver un $\Pi_n(\mathcal{B}) \xrightarrow{\varphi} \mathcal{C}_n$ 2-universel en le sens que pour tout $T$, le foncteur
\[
\underline{\mathrm{Hom}}_{(\mathrm{Top})}(T, \mathcal{B})
\xrightarrow{u \mapsto \varphi \circ \Pi_n(u)}
\tau_1 \underline{\mathrm{Hom}}(\Pi_n(T), \mathcal{C}_n)
\]
soit une équivalence. Pour $n = 1$ on sait que le topos classifiant de $\mathcal{C}_1$ au sens usuel fait l'affaire, mais pour $n = 2$ déja, je doute que ce Pb universel ait une solution. C'est peut-être lié au fait que le « théorème de van Kampen », qu'on peut exprimer en disant que le 2-foncteur $T \mapsto \Pi_1(T)$ des topos localement 1-connexes vers les groupoides transforme (à 1-équivalence près) sommes amalgamées en sommes amalgamées (et plus généralement commute aux 2-limites inductives), n'est sans doute plus
\marginal{cela semble suspect !}
\note{l'étiquette manuscrite de la flèche s'écrit « $u \mapsto \varphi \circ \pi_n(u)$ » ; la note marginale se rapporte aux lignes de $\underline{\mathrm{Hom}}(\Pi_n(T), \mathcal{C}_n)$ et de sa troncature.}

\page{51}
\note{p. 36 de l'auteur ; pagination de la lettre : 4.}
vrai pour le $\Pi_2(T)$. Ainsi, si $T$ est un espace topologique réunion de deux fermés $T_1$ et $T_2$, il n'est sans doute plus vrai que la donnée d'un 1-champs localement constant sur $T$ « équivaut » à la donnée d'un 1-champ localement constant $\mathcal{F}_i$ sur $T_i$ ($i = 1, 2$) et d'une équivalence entre les restrictions de $\mathcal{F}_1$ et $\mathcal{F}_2$ à $T_1 \cap T_2$ (alors que l'énoncé analogue en termes de 0-champs, i.e. de revêtements, est évidemment correct).
\note{la phrase « Ainsi, si $T$ … évidemment correct). » est marquée en marge d'un crochet.}

\marginal{\uncertain{espacer}}
\marginal{(App.~7)}
L'énoncé (B) plus haut rend clair comment expliciter la cohomologie d'un $n$-groupoide $\mathcal{C}_n$. Si $\mathcal{C}_n = \Pi_n(X)$, et si $\mathcal{F}$ est un $(n-1)$-champ localement constant sur $X$, $e^X_{n-1}$ est le $(n-1)$-champ « final », on a une $(n-1)$ équivalence de $(n-1)$-catégories
\[
\Gamma_X(\mathcal{F}) = \mathcal{F}(X) \simeq \underline{\mathrm{Hom}}(e^X_{n-1}, \mathcal{F})
\]
qui montre que le foncteur $\Gamma_X$ « intégration sur $X$ » sur les $(n-1)$-champs localement constants, qui inclut la cohomologie (non commutative) localement constante de $X$ en dim~$\leqslant n-1$, s'interprète en termes de « $(n-1)$-systèmes locaux » sur le groupoide fondamental comme un $\underline{\mathrm{Hom}}(e^{(\mathcal{C}_n)}_{n-1}, \mathcal{F})$ où maintenant $\mathcal{F}$ est interprété comme un $n$-foncteur
\[
\mathcal{C}_n \xrightarrow{\ \mathcal{F}\ } ((n-1)\text{-}\mathrm{Cat})
\]
et $e^{(\mathcal{C}_n)}_{n-1}$ est le $n$-foncteur constant sur $\mathcal{C}_n$, de valeur la $(n-1)$-catégorie finale.

Pour interpréter ceci en notation cohomologique, il faut que j'ajoute, comme « remords » à la lettre précédente, l'interprétation explicite de la cohomologie non commutative sur un topos $X$, en termes d'intégration de $n$-champs sur $X$. Soit $\mathcal{F}$ un $n$-champ de Picard strict sur $X$, il est donc défini par un complexe de cochaines $\mathcal{L}^\bullet$ sur $X$
\[
0 \to \mathcal{L}^0 \to \mathcal{L}^1 \to \mathcal{L}^2 \to \cdots \to \mathcal{L}^n \to 0
\]
concentré en degrés $0 \leqslant i \leqslant n$ (défini à isomorphisme unique près dans la catégorie dérivée de $\mathrm{Ab}(X)$). Ceci dit, les $H^i(X, \mathcal{L}^\bullet)$ (hypercohomologie) \emph{pour $0 \leqslant i \leqslant n$} s'interprètent comme $H^i(X, \mathcal{L}^\bullet) = \pi_{n-i}\Gamma_X(\mathcal{F})$. Si on s'intéresse à tous les $H^i$ (pas seulement pour $i \leqslant n$) on doit, pour tout $N \geqslant n$, regarder $\mathcal{L}^\bullet$ comme un complexe concentré en degrés $0 \leqslant i \leqslant N$

\page{52}
\note{p. 37 de l'auteur ; pagination de la lettre : 5.}
(en prolongeant $\mathcal{L}^\bullet$ par des $0$ à droite). Le $N$-champ de Picard strict correspondant n'est plus $\mathcal{F}$ mais $\mathcal{T}^{N-n}\mathcal{F}$, où $\mathcal{T}$ est le foncteur « espace classifiant », s'interprétant sur les $n$-catégories de Picard strictes comme l'opération consistant à « translater » les $i$-objets en des $(i+1)$-objets, et à rajouter un unique 0-objet ; il se prolonge aux $n$-champs \add{de Picard} « de façon évidente », on espère, de façon à commuter aux opérations d'image inverse de $n$-champs. On aura donc pour $i \leqslant N$
\[
H^i(X, \mathcal{L}^\bullet) = \pi_{N-i}\Gamma_X(\mathcal{T}^{N-n}\mathcal{F}) \qquad (i \leqslant N).
\]
\note{la lettre cursive du foncteur « espace classifiant », tracée à la main, est lue $\mathcal{T}$, comme le $\mathcal{T}_n$ de la p. 50 ; la marge de la p. 53 écrit « $\mathcal{T}\mathcal{F}$ » cerclé. Elle se distingue mal du $\mathcal{C}$ des groupoides.}
Ceci posé, il s'impose, pour tout $n$-champ de Picard strict $\mathcal{F}$ sur $X$, de poser
\[
\boxed{H^i(X, \mathcal{F}) = \pi_{N-i}\Gamma_X(\mathcal{T}^{N-n}\mathcal{F}) \quad \text{si } N \geqslant i, n}
\]
ce qui ne dépend pas du choix de l'entier $N \geqslant \operatorname{Sup}(i, n)$ [\emph{NB} On a un morphisme canonique de $(n-1)$-groupoides,
\[
\mathcal{T}(\Gamma_X\mathcal{F}) \to \Gamma_X(\mathcal{T}\mathcal{F}),
\]
comme le montrent les constructions évidentes en termes de complexes de cochaines, et on voit de même que celui-ci induit des isomorphismes pour les $\pi_i$ pour $1 \leqslant i \leqslant n+1$.].

\emph{NB} On voit en passant que pour un $n$-champ en groupoides $\mathcal{F}$ sur $X$, si on se borne à vouloir définir les $H^i(X, \mathcal{F})$ pour $0 \leqslant i \leqslant n$, on n'a pas besoin sur $\mathcal{F}$ d'une structure de Picard, car il suffit de poser
\[
H^i(X, \mathcal{F}) = \pi_{n-i}(\Gamma_X(\mathcal{F})) \qquad 0 \leqslant i \leqslant n.
\]
Si d'autre part $\mathcal{F}$ est un $n$-Gr-champ (i.e. muni d'une loi de composition $\mathcal{F} \times \mathcal{F} \to \mathcal{F}$ ayant les propriétés formelles d'une loi de groupe) le $(n+1)$ « champ classifiant » est défini, et on peut définir $H^i(X, \mathcal{F})$ pour $i \leqslant n+1$ par
\[
H^i(X, \mathcal{F}) = \pi_{n+1-i}(\Gamma_X(\mathcal{T}\mathcal{F}))
\]
en particulier
\[
H^{n+1}(X, \mathcal{F}) = \pi_0(\Gamma_X(\mathcal{T}\mathcal{F})) = \text{sections de } \mathcal{T}\mathcal{F} \text{ à équivalences près}.
\]
\note{la formule encadrée l'est sur la page ; le « morphisme canonique de $(n-1)$-groupoides » est tel quel.}

\page{53}
\note{p. 38 de l'auteur ; pagination de la lettre : 6.}
\marginal{($\mathcal{T}\mathcal{F}$)}
Mais on ne peut former $\mathcal{T}\mathcal{T}\mathcal{F} = \mathcal{T}^2\mathcal{F}$ et définir $H^{n+2}(X, \mathcal{F})$, semble-t-il \emph{que} si $\mathcal{T}\mathcal{F}$ est lui-même un Gr-$(n+1)$ champ, ce qui ne sera sans doute le cas que si $\mathcal{F}$ est un $n$-champ de Picard strict ….

Venons en maintenant au cas où $\mathcal{F}$ est un $n$-champ \emph{localement constant} sur $X$, donc défini par un $(n+1)$-foncteur
\[
\mathcal{C}_{n+1} \xrightarrow{\ \mathcal{F}\ } (n\text{-}\mathrm{Cat.\ de\ Picard\ strictes}).
\]
Alors, posant pour $0 \leqslant i \leqslant n$
\[
H^i(\mathcal{C}_{n+1}, \mathcal{F}) = \pi_{n-i}(\underline{\mathrm{Hom}}(e^{(\mathcal{C}_{n+1})}_n, \mathcal{F})),
\]
« on a fait ce qu'il fallait » pour que l'on ait un isomorphisme canonique
\[
H^i(\mathcal{C}_{n+1}, \mathcal{F}) \simeq H^i(X, \mathcal{F}),
\]
(valable en fait sans structure de Picard sur $\mathcal{F}$ …). Il s'impose, pour tout $\infty$-groupoide $\mathcal{C}$ et tout $(n+1)$-foncteur
\[
\mathcal{C} \xrightarrow{\ \mathcal{F}\ } (n\text{-}\mathrm{cat.\ de\ Picard\ strictes})
\]
de définir les $H^i(\mathcal{C}, \mathcal{F})$, pour tout $i$, par
\[
H^i(\mathcal{C}, \mathcal{F}) = \pi_{N-i}\underline{\mathrm{Hom}}(e^{\mathcal{C}}_N, \mathcal{T}^{N-n}\mathcal{F})
\]
où on choisit $N \geqslant \operatorname{Sup}(i, n)$. Si $\mathcal{F}$ n'a qu'une Gr-structure (pas nécéssairement de Picard) on peut définir encore les $H^i(\mathcal{C}, \mathcal{F})$ pour $i \leqslant n+1$ par
\[
H^i(\mathcal{C}, \mathcal{F}) = \pi_{n+1-i}\underline{\mathrm{Hom}}(e^{\mathcal{C}}_{n+1}, \mathcal{T}\mathcal{F}).
\]
Dans le cas $\mathcal{C} = \mathcal{C}_{n+1} = \Pi_{n+1}(X)$, il doit être vrai encore (en vertu de (A) plus haut), que cet ensemble est canoniquement isomorphe à $H^{n+1}(X, \mathcal{F}) = \pi_0\Gamma_X(\mathcal{T}\mathcal{F})$ (c'est vrai et bien facile pour $n = 0$). Décrire la flèche canonique entre les deux membres de
\[
H^{n+1}(X, \mathcal{F}) \simeq H^{n+1}(\Pi_{n+1}X, \mathcal{F}) \ ?
\]

Si on veut réexpliciter (A) et (B), en termes du yoga (C), on arrive à la situation suivante :
On a un $(n+1)$-foncteur entre $(n+1)$-groupoides
\[
f_{n+1}\colon \mathcal{C}_{n+1} \to \mathcal{D}_{n+1}
\]
induisant par troncature un $n$-foncteur

\page{54}
\note{p. 39 de l'auteur ; pagination de la lettre : 7.}
\[
f_n\colon \mathcal{C}_n \to \mathcal{D}_n .
\]
On doit avoir alors :

(A$'$) $f_n$ est une $n$-équivalence sss le $n$-foncteur $\varphi \mapsto \varphi \circ f_n$
\[
f_n^*\colon \underline{\mathrm{Hom}}(\mathcal{D}_n, ((n-1)\text{-}\mathrm{Cat})) \to \underline{\mathrm{Hom}}(\mathcal{C}_n, ((n-1)\text{-}\mathrm{Cat}))
\]
allant des $(n-1)$ systèmes locaux sur $\mathcal{D}_n$ (ou $\mathcal{D}_{n+1}$, c'est pareil) vers les $(n-1)$-systèmes locaux sur $\mathcal{C}_n$, est une $n$-équivalence.

(B$'$) $f_n$ est une $n$-équivalence sss pour tout $n$-système local $\mathcal{F}$ sur $\mathcal{D}_{n+1}$,
\[
\mathcal{F}\colon \mathcal{D}_{n+1} \to (n\text{-}\mathrm{Cat}),
\]
le $n$-foncteur induit par $f_{n+1}$
\[
\underbrace{\underline{\mathrm{Hom}}(e^{\mathcal{D}_{n+1}}_n, \mathcal{F})}_{=\, \Gamma_{\mathcal{D}_{n+1}}(\mathcal{F})\ \text{déf}}
\to
\underbrace{\underline{\mathrm{Hom}}(e^{\mathcal{C}_{n+1}}_n, f_{n+1}^*\mathcal{F})}_{=\, \Gamma_{\mathcal{C}_{n+1}}(\mathcal{F})\ \text{déf}}
\]
est une $n$-équivalence.
\note{les accolades et les « $=$ déf » sous les deux membres sont ajoutés à la main ; le second porte $\Gamma_{\mathcal{C}_{n+1}}(\mathcal{F})$, sans $f_{n+1}^*$.}

\marginal{\uncertain{espacer}}
\marginal{(App.~8)}
La construction de la cohomologie d'un topos en termes d'intégration des champs ne fait aucun appel à la notion de complexe de faisceaux abéliens, encore moins à la technique des résolutions injectives. On a l'impression que dans son esprit, via la définition (qui reste à expliciter !) des $n$-champs, elle s'apparenterait plutôt aux calculs « Čechistes » en termes d'hyperrecouvréments. Or ces derniers se décrivent à l'aide d'une petite dose d'algèbre semi-simpliciale. Je ne sais si une théorie des champs et opérations sur iceux pourra s'écrire \emph{sans} jamais utiliser de l'algèbre semi-simpliciale. Si oui, cela ferait essentiellement trois approches distinctes pour construire la cohomologie d'un topos :

a) point de vue des complexes de faisceaux, des résolutions injectives, des catégories dérivées (\emph{algèbre homologique commutative}),

b) point de vue Čechiste ou semi-simplicial (\emph{algèbre homotopique}),

c) point de vue des $n$-champs (algèbre catégorique, ou \emph{algèbre homo-}

\page{55}
\note{p. 39 bis de l'auteur ; pagination de la lettre : 8.}
\emph{logique non commutative}).

Dans a) on « résoud » les coefficients, dans b) on résoud l'espace (ou topos) de base, et dans c) en apparence on ne résoud ni l'un ni l'autre.

Bien cordialement

Alexandre

\section{Lettre à Lawrence Breen, Villecun, 17/19 juillet 1975}

\page{56}
\note{p. 40 de l'auteur. Copie dactylographiée, en anglais, d'une lettre de Grothendieck à Lawrence Breen ; le lieu et la date sont de sa main, ainsi que les corrections. La page suivante porte la pagination de la lettre « -12- » : les pages 2 à 11 de la lettre ne figurent pas dans ce lot.}
Vil\uncertain{le}cun 17/19 July 1975

Dear Larry,

\marginal{(App.~9)}
I am happy to finish by receiving an echo to my long letter and even a beginning to a constructive approach to a theory of the type I envisaged. The construction which you propose for the notion of a non-strict $n$-category, and of the nerve functor, has certainly the merit of existing, and of being a first precise approach, but otherwise can be subject to some evident criticism: it is very technical, unintuitive (yet at the level of 1-cat, etc, and even of 2-cat, everything is so clear -- « you just follow your nose\add{-tip} … »). And finally the absence of a definition of a functor sending (semi-)simplicial sets to $n$-groupoids. This functor corresponds to a geometric intuition so clear that a theory which does not include it seems to me \struck{inadequate.} \add{kind of a joke!} Perhaps in trying to write down (like a sort of list of Christmas presents!) in a complete and explicit enough way the notions which one would like to have at ones disposal, and the relations (functor, equivalence, etc.) which should link them, one would arrive finally at a kind of axiomatic description sufficiently complete which should either give the key to an explicit \emph{ad hoc} construction, or should permit \struck{maybe} \add{at least} to enunciate and prove a theorem of existence and uniqueness\add{(25)} for a theory of the required type.
\note{« theory » : la frappe porte « theorem », corrigé à la main en « theory ». L'appel « (25) » est ajouté à la main.}

Otherwise, not having understood the idea of Segal in your last letter (which I have generously sent to Illusie …), I do not see how you define the Picard $n$-categories -- but this matters little. As far as « strict » Picard $n$-categories are concerned, all I ask of them is that they finally form an $(n+1)$-category $(n+1)$-equivalent to that of \add{chain} complexes \struck{of chains} of length $n$. Agreed? I thank you for having rectified in my mind a big blunder, due to my great ignorance of algebraic topology and homotopy -- I was in fact of the impression that H-spaces satisfying conditions of associativity and commutativity strict enough (say equivalent to an $\Omega^i X$ with $i$ arbitrarily large) correspond to commutative topological groups (inspired by several analogies …). Thus I am entirely in agreement with your observations on p.~5.

\page{57}
\note{pagination de la lettre : « -12- », dont le premier chiffre est surchargé à la main d'un 4 (pagination propre de Grothendieck, 42 ; de même 43 à 45 aux pages suivantes). La suite de la p. 56 n'est donc pas immédiate.}
\marginal{(26)}
On the other hand, I am still intrigued by the following question: is there an analogue of the « tapis » of Dold-Puppe\add{(26)} for semi-simplicial groups (\emph{not necessarily commutative}) and what form should it take? To tell the truth I consider the yoga
\[
(*) \qquad \text{simplicial sets} \longleftrightarrow \infty\text{-groupoids}
\]
as being essentially the ultimate « set theoretic » version of Dold-Puppe, which I would deduce \add{from $(*)$,} by making explicit solely the fact that the abelian groups in ($\infty$-cat) are « nothing else » than the \add{chain} complexes \struck{of chains} in (Ab). One should therefore \add{first} determine what should be the groups in ($\infty$-cat). I can tell you what these are in (1-cat) -- this will be discussed at length in the book of Mme.~Sinh, I think in the chapter « \emph{strict} Gr-categories » (i.e. the isomorphisms of associativity, for unity and inverse $X.X^{-1} \simeq 1$ are \emph{identities}). One can make explicit for example how (\emph{via} the fact that a Gr-category is Gr-equivalent to a strict Gr-category) the calculation with the Gr-categories reduces to a very algebraic calculation with the \emph{strict} Gr-categories, by a kind of « calculus of fractions » (by choice, left or right) of the type which is used in giving the construction of derived categories. In any case, here is the explicit formulation of the structure (groups in (1-cat)) in terms of the theory of groups (1-categories in (Gr))\add{(27)}. The structure is described by a quadruplet $(L_1, L_0, d, \theta)$ with
\marginal{(27)}
\[
L_1 \xrightarrow{d} L_0
\]
a homomorphism of ordinary groups,
\[
\theta\colon L_0 \to \mathrm{Aut}_{\mathrm{Gr}}(L_1)
\]
an operation of $L_0$ on $L_1$, with the following two axioms:

(a) $d$ commutes with the operation of $L_0$, when $L_0$ acts on $L_1$ via $\theta$ and on itself by inner automorphisms:
\[
d(\theta(x_0).x_1) = \mathrm{int}(x_0)\, d(x_1)
\]
(b)
\[
\theta(d(x_1) = \mathrm{int}(x_1) .
\]
These properties imply that $\operatorname{Im} d$ is normal in $L_0$ (hence $\operatorname{Coker} d = \pi_0$ is defined) and $\pi_1 = \operatorname{Ker} d$ is central in $L_1$, and finally that $L_0$
\note{les appels (25), (26), (27), ajoutés à la main dans le texte et répétés en marge, renvoient à des notes qui ne figurent pas dans ce lot. Dans (b), la frappe porte « $\theta(d(x_1) = \mathrm{int}(x_1)$ » ; il manque une parenthèse fermante, laissée telle quelle.}

\page{58}
\note{pagination de la lettre : « -13- », surchargée en 43.}
operates on $L_1$ leaving $\pi_1$ invariant, and it operates \emph{via} $\pi_0$. The principal cohomological invariant of this situation is evidently the Postnikoff-Sinh invariant
\[
\alpha \in H^3(\pi_0, \pi_1) .
\]
I have met these animals -- without even looking for them -- in many situations, which I will not list now (I came across them recently \emph{a propos} the classification of « ordinary » formal groups over a perfect field, in terms of \emph{affine} algebraic groups, and \emph{commutative} formal groups, related by the strict Gr-structures of this type (except that one has to use this formalism in an arbitrary topos (not merely in \struck{\ill{}} \add{(Sets)}) -- to make explicit the yoga that « the transcendent character of a formal group is concentrated essentially in the commutative formal groups », discovered it seems by Dieudonné …). The question which I wish to raise is the generalisation to groups in ($n$-cat), where I expect to find a non-commutative chain complex
\[
L_n \to L_{n-1} \to \cdots \to L_1 \to L_0 \to 1
\]
with supplementary structures doubtless of the type of $\theta$, but what are they? It is understood that the topological significance of \struck{the} \add{such} structures is that they express exactly the « truncated homotopy type in $\dim \leqslant n$ » of topological groups, \struck{and even} \add{or equivalently,} the homotopy type in $\dim \leqslant n+1$ of \struck{connected} \add{pointed connected} topological spaces …). Have you candidates to propose?

\marginal{\uncertain{espacer}}
\marginal{(App.~10)}
Your reflections on biduality and homology, \struck{are attached to} \add{however formal, tie in with} a crowd of developments, of which \add{only} some \struck{only} exist at present, and others would demand \struck{some} considerable \struck{development} \add{work still}. Here are the reminiscences which your naive questions bring to mind:

(A) The formalism of the $Rf_!$, $Rf^!$ (combined with $Rf_*$, $Lf^*$, $\overset{\mathbb{L}}{\otimes}$ and $\mathbb{R}\mathrm{Hom}$, « the six operations ») carries implicitly in itself the definition of homology and the essential identit\struck{ies that underly a}\add{y between} homology and cohomology. One now has this formalism for quasi-coherent sheaves on schemes -- seminar Hartshorne (Springer L.N.~20) -- for the topological spaces and arbitrary sheaves of
\note{le mot rayé avant « (Sets) » est noirci et illisible. Des astérisques de sa main surmontent « biduality » et « homology » ; leur renvoi n'est pas sur la page. Les indices et exposants « ! » sont tracés à la main.}

\page{59}
\note{pagination de la lettre : « -14- », surchargée en 44.}
coefficients -- Verdier, \emph{exposé} Bourbaki (SNLM 300) -- and for the étale cohomology of schemes for « discrete » coefficients (\struck{with} « $\ell$-adic » \add{or} torsion) prime to the residual characteristics \add{(}SGA~5\add{)}, finally, for coherent sheaves on analytic spaces \add{(}Verdier-Ruget\add{)}. (The formalism remains to be develop\struck{p}ed in the crystalline context, and in characteristic $0$ in the context of \add{stratified} modules \struck{stratifications and} \add{with} singularities, à la Déligne, with perhaps -- over the field $\mathbb{C}$ -- the introduction of additional Hodge structures, finally in the context of moti\struck{fs}\add{ves}. I am \struck{persuaded} \add{convinced} that it exists \struck{in some other places -- in all the places, perhaps where} \add{about anywhere -- maybe, wherever} there is a formalism of a cohomological nature.)
\note{« developed » : la frappe porte vraisemblablement « developped », dont la fin est surchargée à la main ; lecture incertaine. La virgule et le tiret de la phrase sur $\mathbb{C}$ sont retouchés à la main.}

Working in étale cohomology on a separated scheme of finite type \struck{on} \add{over} a field $k$, say, with a ring of coefficients $\Lambda$ of torsion prime to the characteristic, the complex of sheaves $f^!(\Lambda_e)$ (where $f\colon X \to \operatorname{Spec} k = e$) plays the rôle of \emph{complex of singular chains on $X$ with coefficients in $\Lambda$}, and $Rf_!(f^!\Lambda_e)$ plays the rôle of a \emph{homology} $H_*(X/e)$, \emph{vis a vis} of course, of coefficients on $e$ which are complexes of $\Lambda$-modules. You can easily justify this assertion with the help of the « global duality theorems », by one or two tricks which I spare you here.
\note{au-dessus du $f$ de « where $f\colon X \to \operatorname{Spec} k$ » un signe $\wedge$ est tracé à la main, sans doute un renvoi d'insertion.}

\emph{REMARKS} (1) There is no \add{need to} truncat\struck{ion}\add{e}, it works in all dimensions.

(2) This is related (at least \struck{from a} \add{as far as the} philosoph\struck{ical point of view}\add{y is concerned}) to the fact that for the various types of coefficients (under conditions of « constructibility ») one has a theorem of « biduality », at least if one allows resolution of singularities (but Deligne has told me I believe that he knows a proof without that), with values in a « dualising complex » $K_e$ (on $e$), $K_X$ (on $X$). If for example $\Lambda$ is « self-dualising » (or Gorenstein) for example $\Lambda = \mathbb{Z}/n\mathbb{Z}$, one can take $K_e = \Lambda$, therefore the dualising complex $K_X = f^!(K_e)$ is nothing else than the « complex of singular chains with coefficients in $\Lambda$ ».

(3) One can do the same thing for coefficients such as $\mathbb{Z}_\ell$ (Jouanolou, \add{thesis --} non published \struck{\uncertain{thesis}}, I fear!).

\page{60}
\note{pagination de la lettre : « -15- », surchargée en 45.}
(4) This works \add{also} for $f\colon X \to S$ finitely presented separated if $f$ has the properties of « cohomological local triviality » (\struck{enhaut local properties} \add{(properties « local upstairs »)}) for example $f$ \emph{smooth}; one finds that $H_*(X/S) = Rf_!f^!(\Lambda_S)$.
\note{dans « $H_*(X/S)$ » les parenthèses sont retracées à la main ; des astérisques de sa main sont placés sous la formule.}

\marginal{\uncertain{espacer}}
(B) Artin-Mazur have studied in a spirit close to yours the \emph{autodualit}\struck{\emph{ies}}\add{y} of the Jacobian\struck{s} of a relative curve $X/S$. It is necessary to ask them for precise results, perhaps it works say if $X/S$ is proper and flat of relative dimension~1 -- in any case it is OK on a discrete valuation ring with smooth \emph{generic} fibre. The special fibre\struck{s} could be very wild. (I have used their results in SGA~7 to prove, in the case of Jacobians, a duality conjecture on the group of connected components associated to the Neron models of abelian varieties dual one to another …)

Towards the end \struck{of decade} of the 50's (beginning of 60's?), when the grand cohomological stuff ($f_!$, $f^!$, étale cohomology, etc.) just came out from darkness, the course given by Serre on the theory of Rosenlicht and Lang on generalised jacobians and the geometric class field theory (see Serre's book) and later the « geometric » theory of \emph{local} class field theory making use of pro-algebraic groups (see his article on this subject), made me reflect on the cohomological formulations of these and other results, which should be of a « geometric » nature, such \struck{as} \add{that} the « arithmetic » results \struck{on the} \add{over an} arbitrary base field (or residual field) $k$ (finite, for example) follow immediately by descent from the « geometric » case of base field $\bar{k}$. I exchanged letters with Serre -- I don't know if I can find copies -- but I recall that \struck{he suggested} \add{I sketched} projects for some ambitious enough theories on generalised residue\struck{ fields}\add{s}, generalised local jacobians, etc., in at least three different directions. But I have never, in spite of numerous attempts, \struck{come to a satisfactory -- position with regard to} \add{succeeded in mobilizing someone for developing one of} these programmes. Here \struck{are several} \add{a few} words on them:

\marginal{\uncertain{espacer}}
\marginal{(App.~11)}
\struck{In the situation} \add{(C) In the situation} where $X$ is of finite type \struck{on a} \add{over a} \emph{field} $k$, construction of a complex of generalised jacobians $J_{*X/k}$ (of length equal to $\dim X$). This is a complex of affine \add{commutative pro-}algebraic \struck{pro}groups on $k$, with the exception of
\note{la lettre continue au-delà de ce lot. Le chiffre de l'appel « (App.~11) » est surchargé et lu avec doute ; les marques marginales lues « espacer » (pp. 45, 46, 51, 54, 58, 60) accompagnent des crochets et des flèches qui délimitent des alinéas, et paraissent des consignes de mise en page.}

\end{document}
